\chapter{FUNDAMENTAÇÃO TEÓRICA}

Este capítulo apresenta a fundamentação teórica que embasa o desenvolvimento deste trabalho, cuja proposta consiste na apresentação de uma plataforma web com o objetivo de centralizar a gestão de clientes, processos, diligências, prazos e indicadores gerenciais de um escritório de advocacia.

\section{TRANSFORMAÇÃO DIGITAL NO SETOR JURÍDICO}

A transformação digital no setor jurídico tem modificado a forma como os escritórios de advocacia e os tribunais desenvolvem suas atividades. Nas últimas décadas, recursos como a virtualização de servidores, a computação em nuvem e as ferramentas colaborativas passaram a ser utilizados nesse ambiente, contribuindo para a digitalização de atividades anteriormente realizadas de forma predominantemente física \cite{pimenta2024}.

Nesse contexto, \citeonline{santos2021} destacam a necessidade de adaptação da área do Direito diante dos avanços tecnológicos e da chegada da Era da Inteligência Artificial. Um dos marcos desse processo foi a implantação do Processo Judicial Eletrônico (PJe), que possibilitou a realização de consultas, protocolos e acompanhamentos processuais por meios digitais \cite{rolimneto2025}. Em 2020, com a pandemia de COVID-19, a transformação digital foi intensificada, levando os tribunais a ampliarem o uso de audiências telepresenciais e ferramentas de videoconferência para dar continuidade às suas atividades.

Além da digitalização dos procedimentos, a incorporação de tecnologias emergentes também tem ampliado as possibilidades de uso de recursos digitais no setor jurídico. Entre essas tecnologias, destaca-se a Inteligência Artificial (IA), que vem sendo utilizada por tribunais brasileiros em atividades como a triagem de petições, o agrupamento de processos repetitivos e o apoio à elaboração de minutas e pareceres. Em mapeamento conduzido pelo Centro de Inovação, Administração e Pesquisa do Judiciário (CIAPJ/FGV Justiça), a Fundação Getulio Vargas identificou 63 iniciativas de IA no Judiciário, das quais 29 estavam em desenvolvimento, 7 eram experiências-piloto e 27 estavam em produção \cite{fgv2020}.

Levantamentos posteriores também buscaram acompanhar a evolução desse cenário. Em pesquisa realizada em parceria entre o Conselho Nacional de Justiça e a Fundação Getulio Vargas \cite{cnjfgv2025}, foram identificados benefícios percebidos no uso da Inteligência Artificial Generativa (IAG), relacionados à eficiência e à qualidade do serviço jurisdicional, além de preocupações quanto à precisão dos resultados, às questões éticas e legais e à utilização informal dessas ferramentas. Nesse contexto, destaca-se o fenômeno conhecido como \textit{shadow AI}, caracterizado pelo uso de ferramentas de IA sem mecanismos institucionais adequados de governança e controle, evidenciando desafios relacionados à segurança e à transparência no uso dessas tecnologias.

Dessa forma, a transformação digital no setor jurídico envolve não apenas a adoção de novas tecnologias, mas também a adaptação dos processos de trabalho e das formas de organização das informações. Nos escritórios de advocacia, esse cenário reforça a importância de uma gestão estruturada, capaz de acompanhar o volume de documentos, atividades e prazos relacionados aos casos atendidos. Nesse sentido, os sistemas de gestão jurídica podem contribuir para centralizar informações, facilitar o acompanhamento das etapas pré-processuais e processuais e aprimorar o controle das atividades do escritório, considerando também a segurança dos dados dos clientes.

\section{SISTEMAS DE GESTÃO E ORGANIZAÇÃO DE ESCRITÓRIOS DE ADVOCACIA}

Nesse contexto de transformação digital, os sistemas de gestão passaram a representar recursos de apoio à organização das atividades dos escritórios de advocacia. Entre os aspectos relevantes dessa organização, destaca-se a gestão documental, que, segundo \cite{aleixo2019gestao}, contribui para a administração eficiente das informações, facilitando sua recuperação e auxiliando a tomada de decisões. Dessa forma, a organização dos documentos constitui um aspecto importante para o funcionamento dos escritórios, especialmente diante da necessidade de consultar informações e acompanhar as atividades relacionadas aos casos atendidos.

A utilização desses sistemas permite centralizar diferentes informações em um único ambiente, contribuindo para uma maior organização das atividades realizadas pelos profissionais do escritório. Além de auxiliar no acompanhamento dos processos, os sistemas de gestão jurídica também podem ser utilizados para automatizar tarefas administrativas e facilitar o controle das atividades realizadas pela equipe.

Pensando nisso, com o avanço da tecnologia, a forma como os escritórios de advocacia realizavam suas rotinas também passou por mudanças. Durante muito tempo, essas atividades eram dominadas pelo uso de documentos físicos, fazendo com que a eficiência dependesse principalmente da organização de cada colaborador. Segundo \citeonline{marques2025}, a gestão jurídica encontra-se atualmente em uma nova etapa, marcada pela presença de tecnologias como a Inteligência Artificial. O autor destaca que essa transformação vai além da automação de tarefas repetitivas, modificando também a forma como os advogados realizam suas atividades, como a pesquisa de informações e a organização do trabalho. Dessa forma, a evolução tecnológica pode proporcionar ganhos de eficiência para os escritórios de advocacia.

Nesse sentido, o gerenciamento de processos é uma atividade que pode ser auxiliada pelos sistemas de gestão. Segundo \citeauthor{abpmp2013} (\citeyear{abpmp2013}), o gerenciamento de processos de negócio é uma forma de organizar o trabalho da equipe, com foco no processo como um todo, desde a análise até o acompanhamento das atividades. Dessa maneira, a utilização de ferramentas tecnológicas pode contribuir para o acompanhamento das atividades de forma mais organizada e para a identificação de possíveis melhorias nos processos realizados.

Pensando no contexto dos escritórios de advocacia, esse gerenciamento pode ser aplicado ao acompanhamento dos processos jurídicos, prazos, documentos e atividades atribuídas à equipe do escritório. A convergência dessas informações em um sistema permite que os colaboradores tenham maior controle sobre as tarefas a serem realizadas, além de favorecer a busca pelo andamento dos processos.

Essa aplicação pode ser observada na dissertação de \citeonline{agostini2010}, que investigou os processos de negócio que compõem a estrutura de escritórios de advocacia sob a perspectiva do BPM e da Administração Legal. A autora identificou processos relacionados à captação e ao cadastro de clientes, à elaboração da primeira peça processual, ao acompanhamento processual e ao encerramento processual. Essa abordagem evidencia que as atividades de um escritório envolvem diferentes etapas que podem ser organizadas e acompanhadas de forma integrada.

Além do gerenciamento dos processos, a colaboração entre os profissionais também é um aspecto importante na organização dos escritórios de advocacia. Segundo \citeonline{aparicio2012}, os sistemas colaborativos são tecnologias projetadas para facilitar o trabalho em grupo, oferecendo mecanismos de controle, coordenação, colaboração e comunicação entre os usuários. O objetivo desses sistemas é reduzir as barreiras impostas pelo espaço físico e pelo tempo. No contexto dos escritórios de advocacia, essas ferramentas podem contribuir para a centralização das informações, o compartilhamento de documentos e o acompanhamento das atividades e prazos de forma integrada.

PRECISAR DAR UMA APROFUNDADA MELHOR NESSA SESSÃO, E TALVEZ CRIAR UMA SESSÃO PARA FALAR SOBRE TECH 
\section{TRABALHOS CORRELATOS}

No âmbito da compreensão sobre a expansão da tecnologia no setor jurídico, um estudo relevante é o artigo ``Legal tech: analytics, inteligência artificial e as novas perspectivas para prática da advocacia'', de \citeonline{andrade2020}. Os autores abordam a utilização de tecnologias como análise de dados e inteligência artificial na atividade jurídica, destacando também soluções tecnológicas empregadas no Poder Judiciário, como o Sistema de Automação da Justiça (SAJ). Embora o estudo tenha um enfoque mais amplo do que a gestão interna de escritórios, contribui para contextualizar a expansão das tecnologias digitais no setor jurídico.

Seguindo essa linha de aplicação tecnológica, \citeonline{castro2020} apresentam o TCC ``Sistema de informação para controle de agendamento de tarefas para escritório de advocacia''. O trabalho apresenta recursos voltados à organização de informações e ao controle das atividades, incluindo agendamentos, dados de clientes e documentos. A proposta se relaciona ao JurisFlow pela preocupação com a centralização das informações e a organização da rotina do escritório. Entretanto, o JurisFlow busca integrar o controle das atividades ao acompanhamento das demandas jurídicas, relacionando diferentes informações e responsáveis em uma mesma estrutura.

Em seguida, o trabalho de \citeonline{rodrigues2021} apresenta o sistema MaisJuridico, no TCC ``Sistema de Agendamentos de Clientes para Advogados'', cuja solução tem como foco o gerenciamento da agenda dos advogados, permitindo registrar consultas, clientes e horários, além de consultar informações relacionadas aos processos. Enquanto essa proposta prioriza a organização dos compromissos, o JurisFlow busca abranger um conjunto mais amplo de atividades, incluindo o acompanhamento das demandas, das diligências e das movimentações registradas pela equipe.

Outro trabalho relacionado à proposta é o artigo ``ADVOCATECH: sistema web para gerenciamento de escritórios de advocacia'', de \citeonline{dias2023}. O estudo apresenta uma aplicação web direcionada à gestão de escritórios, destacando a possibilidade de apoiar a organização das atividades e reduzir o trabalho manual relacionado à localização de documentos e ao acompanhamento de prazos. Sua proximidade com o JurisFlow está no objetivo de utilizar uma plataforma web para apoiar a gestão do escritório. Como característica da proposta, o JurisFlow busca diferenciar as demandas pré-processuais das processuais, permitindo acompanhar atividades anteriores à formalização de um processo judicial.

Já \citeonline{pimenta2024} apresentam o sistema BRSADV, uma plataforma web que oferece recursos relacionados ao atendimento e à atuação jurídica, como consultas virtuais, inserção de documentos e análise prévia do perfil jurídico. O trabalho demonstra a aplicação de tecnologias digitais no apoio às atividades realizadas por profissionais do Direito. Diferentemente dessa abordagem, o JurisFlow concentra-se na organização interna das demandas e no acompanhamento das atividades atribuídas aos responsáveis pelo escritório.

Na sequência, o TCC ``Legal One como ferramenta de agilidade e controle para gestão de escritório de advocacia'', desenvolvido por \citeonline{santos2024}, aborda a utilização do software Legal One como recurso de apoio à gestão jurídica. O estudo destaca aspectos relacionados à automação de atividades, à organização de informações de clientes e processos e ao monitoramento processual. A proposta se aproxima do JurisFlow por tratar da utilização de sistemas para apoiar a gestão e o acompanhamento das atividades de um escritório. O JurisFlow, entretanto, propõe organizar as demandas de acordo com suas etapas, integrando informações e atividades relacionadas a cada caso.

Por fim, o trabalho mais recente analisado é o TCC intitulado ``Juris: desenvolvimento de um sistema web para gerenciamento de processos jurídicos'', desenvolvido por \citeonline{genu2026juris}. O sistema utiliza Spring Boot, React e PostgreSQL para organizar e gerenciar informações jurídicas, tendo como objetivo projetar e implementar um sistema de informação (aplicativo web) focado na construção de funcionalidades para a organização e gerenciamento de processos dentro de um escritório de advocacia. Sua proposta se aproxima do JurisFlow por também utilizar uma aplicação web para gerenciamento de dados relacionados à atividade jurídica.

A análise dos trabalhos correlatos evidencia diferentes aplicações da tecnologia no setor jurídico, abrangendo o gerenciamento de processos, o controle de tarefas, o agendamento de compromissos, o atendimento a clientes e a gestão de escritórios. Embora essas soluções apresentem objetivos e funcionalidades distintos, elas demonstram a utilização de recursos digitais para apoiar a organização das atividades jurídicas. Nesse contexto, o JurisFlow propõe integrar informações de clientes, pastas, casos, processos, diligências, responsáveis e movimentações, contemplando tanto as etapas pré-processuais quanto as processuais. Dessa forma, busca oferecer um ambiente centralizado para o acompanhamento das demandas e das atividades realizadas pela equipe do escritório.

\begin{quadro}[H]
\centering
\caption{Comparação entre os trabalhos correlatos e o JurisFlow}
\label{qua:comparativo}
\renewcommand{\arraystretch}{1.3}
\begin{tabular}{|p{3cm}|p{4.5cm}|p{1.8cm}|p{1.8cm}|p{1.8cm}|}
\hline
\textbf{Trabalho} & \textbf{Foco principal} & \textbf{Controle de prazos} & \textbf{Gestão de demandas} & \textbf{Etapas pré-processuais} \\
\hline
\hline
Andrade et al. (2020) & Legal tech, análise de dados e IA & Não & Indireta & Não \\
\hline
Castro et al. (2020) & Agendamento de tarefas em escritório de advocacia & Não & Parcial & Não \\
\hline
Rodrigues (2021) & Agendamento de clientes e compromissos & Não & Parcial & Não \\
\hline
Dias e Santos (2023) & Gerenciamento de escritórios de advocacia (ADVOCATECH) & Sim & Sim & Não \\
\hline
Pimenta et al. (2024) & Atendimento remoto e atividades jurídicas (BRSADV) & Não & Parcial & Não \\
\hline
Santos (2024) & Gestão de escritório com Legal One & Sim & Sim & Não \\
\hline
Rolim Neto (2025) & Gestão e produção jurídica (Juris Tantum) & Parcial & Parcial & Não \\
\hline
\textbf{JurisFlow (2026)} & \textbf{Gestão de clientes, pastas, casos, processos, diligências, prazos e indicadores} & \textbf{Sim} & \textbf{Sim} & \textbf{Sim} \\
\hline
\end{tabular}
\fonte{Elaborado pelos autores (2026).}
\end{quadro}



O Quadro~\ref{qua:comparativo} sintetiza as principais características dos trabalhos analisados em relação ao JurisFlow. As soluções apresentam diferentes enfoques, incluindo análise de tecnologias jurídicas, agendamento de tarefas, atendimento e gerenciamento de escritórios. Nesse contexto, o JurisFlow propõe integrar informações de clientes, pastas, casos, processos, diligências, responsáveis e movimentações, contemplando tanto as etapas pré-processuais quanto as processuais. Seu diferencial proposto está na organização dessas informações em um único ambiente, permitindo acompanhar as demandas ao longo de suas diferentes etapas.